\section{Entropy compactness and the gadget classification}\label{app:compactness}

This appendix first establishes a general-channel compactness statement. The later steps use the singleton
anchors and untwisted table explicitly. In particular, the finite-mixture conclusion in the first lemma must
not be replaced by a single-factor conclusion for an arbitrary channel.

\subsection{Spectral tightness and finite limiting layers}

\begin{lemma}[Entropy compactness]\label{lem:entropy-compactness}
Let rounds on a fixed $M_d$ have channels $\Psi_j\to\Theta$, bounded source entropy, and production $a_j\to0$.
Then $\Theta$ is a finite convex mixture of finite tracial factorizations.
\end{lemma}
\begin{proof}
Diagonalize source and output bath states in separate decreasing eigenbases, with spectra $p_{ji},q_{jl}$.
Write the system blocks of the unitary as $U_{j,li}$ and put
$w_{j,li}=\Tr(U_{j,li}^*U_{j,li})/d$. Unitarity makes $w_j$ doubly stochastic and $q_j=w_jp_j$.
The probability arrays $w_{j,li}p_{ji}$ and $w_{j,li}q_{jl}$ have relative entropy $a_j$ in bits. Pinsker gives
\[
 \delta_j:=\sum_{li}w_{j,li}|p_{ji}-q_{jl}|\le\sqrt{2\ln2\,a_j}\longrightarrow0.
\]
Sorted matching minimizes absolute-distance transport. Indeed, integrate threshold indicators: any doubly
stochastic coupling mismatches at least the difference between the numbers of eigenvalues above each threshold.
Hence $\sum_i|p_{ji}-q_{ji}|\le\delta_j$.

Sorted eigenvalues obey $p_{ji}\le1/i$. If the entropy bound is $B$, then
$\sum_{i>K}p_{ji}\le{B/\log_2(K+1)}$; the same holds for $q_j$ with a bounded entropy bound.
After adjoining zero coordinates, diagonal subsequence extraction gives pointwise spectral limits, and these
tail bounds upgrade convergence to $\ell^1$. Both limits are the same probability list $p$, with
$S(p)\le\liminf_jS(p_j)\le B$.

All blocks $U_{j,li}$ are contractions. A further diagonal subsequence gives limits $U_{li}$.
If $p_i\ne p_l$, the transport bound forces $U_{li}=0$.
Fix $p_i=\lambda>0$ and choose $K$ with $1/(K+1)\le\lambda/4$. Eventually $p_{ji}\ge\lambda/2$, so
\[
 \sum_{l>K}w_{j,li}\le4\delta_j/\lambda\longrightarrow0.
\]
The analogous row estimate follows from $q_{jl}\ge\lambda/2$. Thus no positive-weight column or row escapes
to infinity. The set $I_\lambda=\{i:p_i=\lambda\}$ is finite, of size $k_\lambda\le1/\lambda$.
In the unitary column identities, terms outside this set vanish: finite unequal-eigenvalue terms vanish
entrywise, tails vanish in squared block norm, and cross-column tails vanish by Cauchy--Schwarz.
Therefore $V_\lambda=[U_{li}]_{l,i\in I_\lambda}$ is a finite square isometry, hence unitary.

Let $\Theta_\lambda$ be the channel of this finite tracial factorization. Splitting the source at a finite index cutoff makes its
tail channel mass uniformly small. On each retained positive-weight column the preceding no-escape estimate
permits taking the blockwise limit. This proves
\[
 \Theta=\sum_{\lambda>0}k_\lambda\lambda\,\Theta_\lambda,
 \qquad\sum_{\lambda>0}k_\lambda\lambda=1.
\]
It remains to reduce a countable barycenter to a finite one. In a finite-dimensional real space, the
barycenter of bounded atoms with strictly positive weights lies in their convex hull. To see this, if the
barycenter lies on a relative boundary of the closed convex hull, a supporting functional has its maximum
at the barycenter. Every positive-weight atom must attain that maximum; restrict to its lower-dimensional
affine hyperplane. Repeat until the barycenter is in relative interior. A small simplex around it can then
be approximated by vertices in the convex hull, still containing the barycenter. Carath\'eodory elimination
leaves at most $d^4+1$ of the original Choi atoms. Their dimensions are finite, proving the lemma.
\end{proof}

\subsection{Exact support, rational weights and separation}

\begin{lemma}[Support-preserving copies]\label{lem:support-copies}
Let $\Theta=\sum_ip_i\Theta_i$ be a finite mixture of finite tracial factorizations of bath dimensions $m_i\le m$.
For $0<\gamma<1$, it has an approximant $\Gamma$ with a finite tracial factorization of bath dimension at most
$\lceil(d^4+1)m/\gamma\rceil$, half-diamond error less than $\gamma$, and
$\cJ(\Gamma)\le\cJ(\Theta)/(1-\gamma)$.
In particular, a common exact Choi support is preserved.
\end{lemma}
\begin{proof}
Reduce to $L\le d^4+1$ original atoms. Put $N=\lceil(d^4+1)m/\gamma\rceil$,
$c_i=\lfloor Np_i/m_i\rfloor$ and $N'=\sum_ic_im_i$. The missing fraction
$\delta=(N-N')/N$ is less than $\gamma$, so $N'>0$.
Take the block-diagonal unitary with $c_i$ copies of each factor, on a bath of dimension $N'$, as in the proof
of~\cite[Theorem~3.6]{HaagerupMusat2015}.
Its channel $\Gamma$ obeys $\Theta=(1-\delta)\Gamma+\delta\Theta_{\rm rem}$ for a channel
$\Theta_{\rm rem}$ (omit it if $\delta=0$). This proves both bounds. There is no filler channel.
\end{proof}

A finite tracial factorization with exact target support has a unitary expansion $\sum_gA_g\otimes B_g$
(compare~\cite[Theorem~2.2]{HaagerupMusat2011}):
each bath-entry Kraus vector must belong to the target Kraus span by Choi positivity.
The anchors make all square $B_g$ unitary, and the triangle columns give $B_1^*B_y=B_x^*B_z$.
After gauging $V_g=B_1^*B_g$, these are exactly the relations of $P_\Phi$.
The two-anchor reference input gives
\begin{equation}\label{eq:anchor-gap-classification}
 |\tau(V_g^*V_h)|\le2\ddia(\Theta,\Phi),\qquad g\ne h.
\end{equation}
If no finite representation separates the whole finite label set, some fixed pair is equal in every finite
representation. Otherwise a finite direct sum of pair-separating representations separates all pairs.
Thus every exact-support finite tracial mixture is at diamond distance at least $1/2$ from $\Phi$.

\begin{proof}[Bounded branch of Theorem~\ref{thm:trichotomy}]
Suppose $B=\sup_nL^*_\Phi(n)<\infty$. Near-minimizers have bounded $S$, $a\to0$ and $\ell\to0$.
Their channels have a convergent subsequence with limit $\Theta$ at distance at most $1/16$ from $\Phi$,
with exact target support. The compactness proof gives spectral layers with weights
$w_\lambda=k_\lambda\lambda$ and
\[
 S(p)=H(w)+\sum_\lambda w_\lambda\log k_\lambda\le B.
\]
All layers have exact support. Discard those with $\log k_\lambda>16(B+1)$; their mass is below $1/16$.
The remaining normalized barycenter has bath cap $2^{16(B+1)}$. Even if it has infinitely many layers,
the preceding barycenter argument leaves at most $d^4+1$ of them, without changing the channel.
Lemma~\ref{lem:support-copies} at $\gamma=1/32$ gives one exact-support factor of dimension
at most $33(d^4+1)2^{16(B+1)}$, at distance less than $5/32$ from $\Phi$.
Equation~\eqref{eq:anchor-gap-classification} gives a separating representation with trace gap below $5/16$.
This proves the displayed bound on $R_{\rm sep}$.

Conversely, tensor a representation with gap at most $1/2$ to degree
$t=\lceil\log_2(8r(r-1))\rceil$. The exact unitary table has zero production and leakage and
half-diamond error at most $r(r-1)2^{-t}/2\le1/16$.
It gives $B\le tR_{\rm sep}$. A representation merely separating the labels acquires a strict trace gap by
adding a trivial summand, then a finite tensor power makes the gap at most $1/2$.
Robust rank streaming (Theorem~\ref{thm:robust}(c)) gives logarithmic service memory.

Finally the matrix image of $P_\Phi$ is finitely generated. Malcev's residual-finiteness theorem~\cite{Nica2013}
gives a finite quotient separating each of the finitely many label pairs; take their finite product.
Its regular representation has trace overlaps exactly $\delta_{gh}$ and obeys all relations.
It therefore gives an exact finite tracial factorization of $\Phi$.
The converse follows from the exact-support expansion above. The same fixed-pair argument proves that an
exact finite tracial mixture for this gadget also implies a single exact finite tracial factorization.
\end{proof}

\subsection{Ultrapower separation and closure}

\begin{lemma}[Ultrapower label criterion]\label{lem:ultrapower-labels}
A canonical untwisted anchored gadget is in $\overline{\mathcal F}$
if and only if some homomorphism
$P_\Phi\to U(\mathcal R^\omega)$ separates its named labels.
If there is no such homomorphism, some fixed distinct pair is
identified by every such homomorphism.
\end{lemma}
\begin{proof}
For a channel in the closure, Haagerup--Musat's factorization
in $\mathcal R^\omega$ gives a unitary whose coefficient expansion
has exact canonical support
\cite[Theorem~3.6]{HaagerupMusat2015}; the expansion is that
of~\cite[Theorem~2.2]{HaagerupMusat2011} (compare~\cite[Proposition~3.9]{Musat2025}).
Faithfulness of the trace makes every orthogonal Kraus coefficient
zero. The singleton anchors make the remaining coefficients
unitary; triangle column orthogonality gives the relations of
$P_\Phi$ after gauging by the identity coefficient. Equality of
the target Choi matrix makes distinct-label trace overlaps zero.
This separates the labels.

Conversely, let $\pi$ separate the labels. Work throughout with
finite matrix approximants of the finitely many images $\pi(v_g)$,
so that every construction below stays inside matrix algebras with
normalized trace. Replace $\pi$ by $\pi\oplus1$, with equal tracial
weights on the two summands. For $g\ne h$ the overlap in
$\mathcal R^\omega$ becomes $(1+\tau(\pi(v_g)^*\pi(v_h)))/2$, of modulus
strictly below one: there the trace is faithful, so for a unitary $u$,
$|\tau(u)|\le1$ with equality only when $u$ is a scalar unitary
$e^{i\theta}I$; then the overlap is $(1+e^{i\theta})/2$, and $\theta=0$
is excluded because $\pi$ separates the labels. So the finitely many
overlaps have modulus at most $1-\delta$ for some $\delta>0$, and along
$\omega$ those of the matrix approximants have modulus at most
$1-\delta/2$.
Tensor powers, or tensor-conjugate powers (the tensor-power trick
of~\cite[Lemma~2.2]{SlofstraVidick2018}), drive these finitely
many overlaps to zero while preserving every relation up to defects
that vanish along the ultrafilter. The resulting finite matrix
tuples have vanishing triangle defects. Form the actual physical gadget
blocks with lower-right entry $-V_xV_y$. Their slot Choi matrices
approach the canonical table, because equal-label slot defects
and distinct-label overlaps tend to zero. This yields finite tracial
factorizations converging to $\Phi$; the same amplification is quantified
in Appendix~\ref{app:hyperlinear-upper}.

Finally, if each pair were separated by some homomorphism, the
finite direct sum of those maps would separate every pair and
embed in $\mathcal R^\omega$. Its contrapositive gives the fixed
pair assertion. This argument concerns finite label separation,
rather than hyperlinearity of all of $P_\Phi$.
\end{proof}

\subsection{Closure-only and nonclosure branches}

At any fixed entropy bound $B$ and coarse distance $u<1/2$, absence of finite label separation implies
\begin{equation}\label{eq:fixed-entropy-gap}
 \inf\{a+\ell:S\le B,\ \ddia\le u\}>0.
\end{equation}
Otherwise compactness would give an exact-support finite mixture at distance below $1/2$, contradicting
the anchor obstruction. In particular $L^*_\Phi(n)\to\infty$ on the nonseparating branch.
If $\Phi\in\overline{\mathcal F}$, positive-Gram repair and scalar completion
(Lemma~\ref{lem:generic-repair}) give exact-target flat rounds with finite $S_\delta$ and $a_\delta\to0$.
Then $L^*_\Phi(n)\le S_\delta+na_\delta$; take $n\to\infty$ and then $\delta\to0$ to obtain $o(n)$.
The closure service upper is Theorem~\ref{thm:closure-law}(a).
If service memory had a bounded subsequence, the precise extracted production bound of
Theorem~\ref{thm:lower} would give $a=O((1+\log n)/n)$ at bounded $S$, contradicting
\eqref{eq:fixed-entropy-gap}. This proves the middle branch.

For the last branch we need the following qualitative amplification, specific to the gadget.
\begin{lemma}[Support amplification]\label{lem:support-amplification}
If $\Theta\in\overline{\mathcal F}$ has exact target support and $\ddia(\Theta,\Phi)<1/2$, then
$\Phi\in\overline{\mathcal F}$.
\end{lemma}
\begin{proof}
Take channels $\Psi_j\to\Theta$ with finite tracial factorizations, with unitaries $U_j=\sum_gA_g\otimes X_{jg}+L_j$,
$\|L_j\|_2^2=\ell_j\to0$ in normalized joint trace. Polar extensions of their singleton anchor blocks
give square unitaries $R_{jg}$ with
$\|X_{jg}-R_{jg}\|_2\le2\sqrt{d\ell_j}$.
The table $W_j=\sum_gA_g\otimes R_{jg}$ has norm at most $\sqrt2$ and
$\|U_j-W_j\|_2\le(1+2\sqrt d)\sqrt{\ell_j}$.
Its column defects therefore tend to zero. Gauging by $R_{j1}^*$ gives unitaries $V_{jg}$,
$V_{j1}=I$, with triangle error at most
$\kappa_j=2\sqrt d(1+\sqrt2)(1+2\sqrt d)\sqrt{\ell_j}$.
Anchor testing and polar error give trace overlaps at most
$2\ddia(\Psi_j,\Phi)+2\sqrt{d\ell_j}$, eventually bounded by $c_0<1$.

Construct actual unitary gadget blocks with lower-right entry $-V_{jx}V_{jy}$, not $-V_{jz}$.
Every slot response $Q_a$ is within $\kappa_j$ of its canonical label response. Equal-label slots therefore
have $1-|\tau(Q_a^*Q_b)|^2\le4\kappa_j^2$; different-label slots eventually have overlap at most $c_1<1$.
Replace all unitary responses by $(V\otimes\overline V)^{\otimes k}$.
This multiplicative map preserves actual unitary blocks; scalar slot signs remain outside the map.
Slot correlations become $|\tau(Q_a^*Q_b)|^{2k}$.
There are at most $2d$ slots and their squared scalar coefficients sum to $d$.
The Choi error is thus at most $d(4k\kappa_j^2+c_1^{2k})$.
Choose $k$ first and $j$ second to make this arbitrarily small, proving closure membership.
\end{proof}

For nonclosure $\Phi$, the compact set
$\{\Theta\in\overline{\mathcal F}:\ddia(\Theta,\Phi)\le1/8\}$ has a strictly positive minimum leakage
$\eta$ by this lemma; if the set is empty, use $\eta=1$.
Put $\eta'=\min(\eta,1)$ and
$a_* =\min\{1/(2048d^2\ln2),\eta'^2/(32\ln2)\}$.
A qualifying round with $a\le a_*$ tracializes to a channel in closure at Choi distance
$h\le\sqrt{2\ln2 a_*}\le\min(1/(32d),\eta'/4)$.
Its target diamond distance is at most $1/16+dh\le3/32<1/8$ and its leakage is at most $2/(15n)+h$,
below $\eta$ for sufficiently large $n$. This is impossible. Hence
$L^*_\Phi(n)\ge na_*$ eventually. The pure Stinespring round gives $L^*_\Phi(n)\le n\log r$.
Theorem~\ref{thm:lower} gives the linear rank-rate memory lower at every $\epsilon\le1/16$;
the exact independent-cell service gives the linear upper. This completes Theorem~\ref{thm:trichotomy}.

\subsection{Regime (iii) and non-hyperlinear groups}\label{app:regime-iii}

A countable group is hyperlinear exactly when it is isomorphic to a subgroup of
$U(\mathcal R^\omega)$~\cite[Remark~8.4 and Theorem~8.5]{Pestov2008}.

\begin{proof}[Proof of Proposition~\ref{prop:regime-iii}]
(a)$\Rightarrow$(b). By Theorem~\ref{thm:trichotomy}(iii), distinct labels $g\ne h$ are identified by
every homomorphism $P_\Phi\to U(\mathcal R^\omega)$, and $P_\Phi$ is finitely presented. Put
$\zeta=v_g^{-1}v_h$. Then $q(\zeta)=g^{-1}h\ne1$, so $\zeta\ne1$, and every such homomorphism kills $\zeta$.

(b)$\Rightarrow$(a). Let $P=\langle S\mid R_0\rangle$ with $S$ and $R_0$ finite. Drop empty relators and
replace each relator $\sigma=s^{\pm1}$ of length one by its conjugate $s\sigma s^{-1}$; the normal
closure is unchanged. Choose a word $w$ of length at least two representing $\zeta$, appending $ss^{-1}$
if necessary. Let $\Phi$ be the canonical gadget of $G=P$, the generating set $S$ and the list
$R=R_0\cup\{ww^{-1}\}$. Every word in $R$ has length at least two and equals $1$ in $P$. The tagged
prefix of $ww^{-1}$ at position $|w|$ names $\zeta$, so $\zeta$ and $1$ are distinct labels.

In $P_\Phi$ the triangle $(s,s^{-1},\varnothing)$ gives $v_{[s^{-1}]}=v_{[s]}^{-1}$. For a word
$\sigma_1\cdots\sigma_m\in R$, the triangles $(p_{i,j-1},\sigma_j,p_{i,j})$ give
$v_{[p_{i,j}]}=v_{[\sigma_1]}\cdots v_{[\sigma_j]}$ by induction on $j$; at $j=m$ the prefix is empty, so
$v_{[\sigma_1]}\cdots v_{[\sigma_m]}=1$. Hence $s\mapsto v_{[s]}$ defines a homomorphism
$\theta:P\to P_\Phi$, and $\theta(\zeta)=v_{[w]}=v_\zeta$. For every homomorphism
$\pi:P_\Phi\to U(\mathcal R^\omega)$, the composite $\pi\theta$ kills $\zeta$, so $\pi(v_\zeta)=1=\pi(v_1)$.
Theorem~\ref{thm:trichotomy} places $\Phi$ in regime (iii).

(b)$\Rightarrow$(c). A homomorphism that kills $\zeta\ne1$ is not injective, so $P$ is not hyperlinear.

(c)$\Rightarrow$(b). Let $P$ be finitely presented and suppose that every $\zeta\ne1$ survives some
homomorphism $\pi_\zeta:P\to U(\mathcal R^\omega)$. Enumerate $P\setminus\{1\}$ as
$\zeta_1,\zeta_2,\ldots$ and choose orthogonal projections $p_k\in\mathcal R$ with $\tau(p_k)=2^{-k}$.
Each corner $p_k\mathcal Rp_k$ is a hyperfinite $\mathrm{II}_1$ factor, so there is a $*$-isomorphism
$\theta_k:\mathcal R\to p_k\mathcal Rp_k$. It scales $\nrm\cdot_2$ by $2^{-k/2}$, so applying it to
representing sequences gives an injective $*$-homomorphism
$\theta_k^\omega:\mathcal R^\omega\to p_k\mathcal R^\omega p_k$ with $\theta_k^\omega(1)=p_k$. The
orthogonal sums $\pi(x)=\sum_k\theta_k^\omega(\pi_{\zeta_k}(x))$ converge in $\nrm\cdot_2$ on bounded sets
and define a homomorphism $\pi:P\to U(\mathcal R^\omega)$. Since
$p_k\pi(\zeta_k)p_k=\theta_k^\omega(\pi_{\zeta_k}(\zeta_k))\ne p_k$, the map $\pi$ is injective, so $P$
is hyperlinear. Contrapositively, a finitely presented group that is not hyperlinear has an element
$\zeta\ne1$ killed by every homomorphism into $U(\mathcal R^\omega)$.
\end{proof}
